\section*{第 \ref{ch:b3:statistical-thermo-applied} 章——统计热力学应用}

\begin{solution}{exo:b3:statistical-thermo-applied:1}
$\frac32R + R + (3 \times 3 - 5)R = 6.5R = \qty{54.0}{J.K^{-1}.mol^{-1}}$。室温下平动
和转动可视为经典的，两个伸缩振动被冻结，二重简并的
弯曲振动（即最低的振动）部分被激发。
\end{solution}

\begin{solution}{exo:b3:statistical-thermo-applied:2}
$x = 1$: $\eu/(\eu - 1)^2 = 0.921$. $x = 10$: $100\eu^{10}/(\eu^{10} - 1)^2 \approx 100\eu^{-10} = \num{4.5e-3}$.
\end{solution}

\begin{solution}{exo:b3:statistical-thermo-applied:3}
\qty{20}{K}: $x_{\mathrm{para}} = 1/(1 + 3 \times 3\eu^{-2 \times 85.35/20}) = 1/(1 + 9 \times \num{1.96e-4}) = 0.998$.
\qty{77}{K}：偶数项之和 $1 + 5\eu^{-6.651} = 1.0065$；奇数项之和 $3(3\eu^{-2.217} + 7\eu^{-13.30}) = 0.9806$；
$x_{\mathrm{para}} = 1.0065/1.9871 = 0.51$.
\end{solution}

\begin{solution}{exo:b3:statistical-thermo-applied:4}
对于 $I = 1$，有 $3 \times 2 = 6$ 个对称、$3 \times 1 = 3$ 个反对称的自旋态；
氘核是玻色子，故总波函数对称：对称的自旋态
配以偶数 $J$（正型，权重为 6），反对称自旋态配以奇数 $J$（仲型，权重为
3）。含 $J = 0$ 的正型氘在 0 K 时是稳定的；室温下
正型：仲型 = 2:1。
\end{solution}

\begin{solution}{exo:b3:statistical-thermo-applied:5}
$\Delta_rE_0 = 2 \times 1883.7 - 2170.3 - 1542.3 = \qty{54.8}{cm^{-1}}$；$\eu^{-54.8/207.2} = 0.768$，且
$4 \times 0.768 = 3.07$。完整值 3.26 高出 6\,\%：平动、转动和
振动因子只有在经典极限下才严格抵消，而
这些轻分子的转动在 \qty{298}{K} 时仍部分量子化。
\end{solution}

\begin{solution}{exo:b3:statistical-thermo-applied:6}
\ce{N2O}：两种取向，$R\ln2 = \qty{5.76}{J.K^{-1}.mol^{-1}}$。\ce{CH3D}：四种，$R\ln4 =
\qty{11.53}{J.K^{-1}.mol^{-1}}$（为上限，仅当排列完全随机时才能达到）。
\end{solution}

\begin{solution}{exo:b3:statistical-thermo-applied:7}
$x = 797.6/300 = 2.659$：爱因斯坦函数 $x^2\eu^{-x}/(1 - \eu^{-x})^2 = 0.572$；$C_V/R = 3.072$，
$C_p = 4.072R = \qty{33.86}{J.K^{-1}.mol^{-1}}$，比表中值低 0.4\,\%（非简谐效应）。
\end{solution}

\begin{solution}{exo:b3:statistical-thermo-applied:8}
$x = 1.029$：$x^2\eu^{-x}/(1 - \eu^{-x})^2 = 0.916$，即 \qty{7.62}{J.K^{-1}.mol^{-1}}，为 $R$ 的 92\,\%：碘的
振动在室温下几乎是经典的。
\end{solution}

\begin{solution}{exo:b3:statistical-thermo-applied:9}
B 中 \ce{^{18}O} 更富集。$\alpha_{\mathrm{A/B}} = (1 - 0.0100)/(1 + 0.0020) = 0.98802$。
\end{solution}

\begin{solution}{exo:b3:statistical-thermo-applied:10}
$\Delta\mathrm{ZPE} = \frac12(1 - 1/\sqrt2)(2000 - 3000) = \qty{-146}{cm^{-1}}$; $K = \eu^{146.4/207.2} = 2.0$.
氘优先进入较硬的键 X，在该键中
其零点能的节省最大。
\end{solution}

\begin{solution}{exo:b3:statistical-thermo-applied:11}
$\Delta_rH^\circ = R\ln10\,(-1.766 + 3.392)/(1/900 - 1/1100) = 19.145 \times 1.626/\num{2.020e-4} =
\qty{154}{kJ/mol}$。两个原子共携带 $2 \times \frac52RT$ 的焓，分子则为 $\frac52RT + RT$
（转动）加上其振动能 $R\theta_{\mathrm V}/(\eu^{\theta_{\mathrm V}/T} - 1)$，略小于 $RT$：
两者之差在 \qty{1000}{K} 时约为 \qty{5}{kJ/mol}，计入 $D_0$。
\end{solution}

\begin{solution}{exo:b3:statistical-thermo-applied:12}
能级 0（$g = 1$）和 $6hcB$（$g = 5$）：$q = 1 + 5\eu^{-x}$，$\langle\varepsilon\rangle/kT = 5x\eu^{-x}/q$，$\langle\varepsilon^2\rangle/(kT)^2 =
5x^2\eu^{-x}/q$，且 $C/R = \langle\varepsilon^2\rangle/(kT)^2 - (\langle\varepsilon\rangle/kT)^2 = 5x^2\eu^{-x}/(1 + 5\eu^{-x})^2$。当
\qty{100}{K}时，$x = 5.121$：$C/R = 0.783/1.061 = 0.738$。
\end{solution}

\begin{solution}{pb:b3:statistical-thermo-applied:1}
\textbf{1.} $2 \times 2 = 4$：三个对称态（三重态），一个反对称态（单重态）。
\textbf{2.} 质子是费米子，因此交换两个质子时总波函数改变符号；
转动函数给出 $(-1)^J$：偶数 $J$ 对应
反对称的单重态（仲型），奇数 $J$ 对应对称的三重态（正型）。
\textbf{3.} 偶数能级权重为 1，奇数能级权重为 3。
\textbf{4.} 在 \qty{300}{K}（$T = 3.5\,\theta_{\mathrm R}$）时，偶数与奇数的转动求和几乎相等；
分别以 1 和 3 加权后得 1:3（仲型比例为 0.2507）。
\textbf{5.} $F(1) = 2 \times 59.322 - 4 \times 0.0471 = \qty{118.46}{cm^{-1}} = \qty{1.417}{kJ/mol}$.
\textbf{6.} $1/(1 + 9\eu^{-170.4/20.37}) = 1/(1 + 9 \times \num{2.3e-4}) = 0.998$.
\textbf{7.} 只考虑 $J = 0$ 和 1，仲型占一半意味着 $9\eu^{-2\theta_{\mathrm R}/T} = 1$，$T = 2\theta_{\mathrm R}/\ln9 =
0.91\,\theta_{\mathrm R} = \qty{78}{K}$：竞争发生在 $J = 1$ 的权重 9 与其
玻尔兹曼因子之间。
\textbf{8.} 0.75。无催化剂时转化需数天，液化
只需数小时。
\textbf{9.} $(0.75 - 0.002) \times 1.417 = \qty{1.060}{kJ/mol}$.
\textbf{10.} $1.060/\num{2.016e-3} = \qty{526}{kJ/kg}$.
\textbf{11.} $0.905/\num{2.016e-3} = \qty{449}{kJ/kg}$.
\textbf{12.} 转化热是汽化焓的 1.17 倍。
\textbf{13.} 它产生于液体内部，而非经由器壁传入。
\textbf{14.} $x_{\mathrm o}(t) = x_{\mathrm{eq}} + (0.75 - x_{\mathrm{eq}})\eu^{-kt}$，其中 $x_{\mathrm{eq}} = 0.002$。
\textbf{15.} $0.002 + 0.748\eu^{-0.274} = 0.571$.
\textbf{16.} $(0.75 - 0.571) \times 1.417 = \qty{0.254}{kJ/mol}$.
\textbf{17.} $0.254/0.905 = 0.28$：一天内储罐中的 28\,\% 被耗去。
\textbf{18.} 经过 \qty{168}{h} 后，$x_{\mathrm o} = 0.112$，热量为 \qty{0.904}{kJ/mol}，与
汽化焓相等：按此模型，储罐已经空了。这一估计高估了损失，
因为蒸发的分子会带走尚未转化的正氢；
实际损失较小，但仍然很大。
\textbf{19.} $t_{1/2} = \ln2/0.0114 = \qty{61}{h}$.
\textbf{20.} 在液化器中，于逐级降温的催化剂床上进行，以使
制冷机在液体储存前即移走转化热。
\textbf{21.} 是：\qty{300}{K} 时平衡仲型比例为 0.2507。
\textbf{22.} 在普通氢中，两种形态不能相互转化，因此各自
的转动被冻结在各自的能级阶梯上；在平衡氢中，
所供热量的一部分使仲型转化为正型。
\textbf{23.} 在 \qty{50}{K} 附近约为 3.57：除激发转动外，热量还使
分子从 $J = 0$ 到 $J = 1$（高出 \qty{118}{cm^{-1}}），这是一个焓为正的
反应。
\textbf{24.} 1.50：两种形态的转动均被冻结（仲型在 $J = 0$，正型在 $J = 1$）。
\textbf{25.} \textbf{转化热 / 汽化焓 $\approx 1.2$，适用于普通氢
}：仅其转化即可蒸发整个储罐，这也是
氢在储存前需转化为仲型的原因。
\end{solution}
